\chapter{Bringing λProlog to \texorpdfstring{\texttt{(aux kanren micro)}}{(aux kanren micro)}: the design}\label{ch:design}

% ------------------------------------------------------------------------------------------------
\section*{What and why}

\Cref{ch:micro} ended with a precise list of what an importing module cannot obtain from \micro{}
today: a named place in the state for its own data, a call after every binding (to check the scope of
eigenvariables eagerly and to wake postponed equations), and a say in reification (to print new
constraints and to hide what must not escape). The three tricks of \cref{sec:micro-noedit} show that
\rulename{augment} and \rulename{generic} \emph{fit} the shape of a μKanren goal; they do not make them sound under every
goal of the engine, nor cheap.

This chapter turns that observation into a design. It proceeds in four steps.
\begin{enumerate}
\item \textbf{Requirements} (\cref{sec:design-req}). The nine acceptance items A1--A9 against which the
  proposal was built and tested, stated as properties of a run of \scm{°->list}.
\item \textbf{The design space} (\cref{sec:design-space}). The four meanings of unification on λ-terms
  compared in \cref{sec:lkanren-compare}, and why higher-order pattern unification with universe levels is
  the one retained.
\item \textbf{Three prototypes} (\cref{sec:design-prototypes}). Three complete implementations of the
  same requirements, which differ in where the code sits with respect to \micro: a layer that does not
  touch the engine, a small generic protocol in the engine, and a native engine. They were compared by
  three reviews (engineering, fidelity to the book, pedagogy); the protocol won.
\item \textbf{The extension protocol} (\cref{sec:design-protocol}). The edit to \micro{} in full: the
  seventh state field, the registry, every hook with its contract and call site, the two new primitives,
  the one generic change to \scm{absent°}, and the two invariants a client must keep.
\end{enumerate}
The chapter closes with what the edit costs a first-order program (\cref{sec:design-cost}) and with
the history of the review rounds that hardened the proposal (\cref{sec:design-reviews}).
The λ layer that sits on the protocol, \scm{(aux kanren lambda)}, with its unifier, its \scm{∀°} and
\scm{⊃°} and its treatment of micro's stores, is the subject of \cref{ch:module}; the book's examples
run on it are in \cref{ch:examples}; the plan for merging is in \cref{ch:roadmap}. In the repository the
proposal lives under \texttt{doc/hol-kanren/prototype/} (the module, \texttt{micro.diff},
\texttt{show.diff}, the design notes \texttt{DESIGN.md}, the tests, the examples, the benchmarks and the
scripts that build and run them, \cref{sec:roadmap-run}). After merging, the module and its suites
become \texttt{src/}\allowbreak\texttt{aux.kanren.lambda.scm} and
\texttt{src/}\allowbreak\texttt{test/}\allowbreak\texttt{microkanren-lambda*.scm}.

Every Scheme snippet of this chapter with an answer was run against the proposal (the edited micro and
\scm{(aux kanren lambda)}) unless it says otherwise, and its output is pasted verbatim. Snippets of the
form \scm{(show "label" query)} are quoted from the proposal's \texttt{examples.scm}, where \scm{show}
prints the label as a comment and then the answers; their outputs are quoted from
\texttt{examples.out}.

% ================================================================================================
\section{Requirements}\label{sec:design-req}

The proposal was written against nine acceptance items, A1 to A9. Each one is a property of the
answers that \scm{°->list} returns for a class of goals, illustrated by the book's own examples with the
book's own answers. This section states each item in a form precise enough to be tested, then gives
its intuition and one worked example. The general statements are \emph{tested}, by 83 hand-written cases
with 320 assertions, not proved.

\subsection{Setting and notation}\label{sec:design-notation}

\paragraph{Goals and terms.} The goal language is that of \micro{} (\scm{=°}, \scm{≠°}, \scm{absent°},
\scm{symbol°}, \scm{number°}, \scm{fresh°}, \scm{and°}, \scm{or°}, \scm{cond°}, relations) extended with
three forms: \scm{(∀° (x) G)}, the universal goal \lp{pi x\ G}; \scm{(⊃° (D ...) G)}, the implication
\lp{D => G}; and relations declared with \scm{define-dynamic-relation}, the only ones that take
hypotheses. The term language is micro's (Scheme data and logic variables) extended with abstractions
\scm{(ƛ (x) t)}, applications \scm{(@ t u ...)} and eigenvariables. We write $\mathrm{Ans}(G)$ for the list
returned by \scm{(°->list #f G)}, and say that $G$ \emph{fails} when $\mathrm{Ans}(G)$ is empty.

\paragraph{Reference semantics.} The specification is the goal-directed search of the book for
hereditary Harrop formulas, that is, the six rules \rulename{and}, \rulename{or}, \rulename{instan}, \rulename{generic}, \rulename{augment} and \rulename{backchain}
(\cref{fig:rules-right}; \book{77}, \book{82}, \book{124}; see \cref{sec:hh-rules} and \cref{app:rules}). Uniform proofs are complete for intuitionistic provability of
these formulas (\cref{thm:hh-intuitionistic}), so implementing these rules faithfully \emph{is} the
semantics. Where λProlog says nothing (disequality, \emph{absento}, types), the reading adopted extends
micro's own semantics over λProlog's term model (item A8).

\begin{definition}[Universe levels]\label{def:design-levels}
A state carries a \emph{depth} $d\in\mathbb N$, the number of \scm{∀°} goals whose body is running.
Every logic variable $X$ has a \emph{level} $\ell(X)$, the depth at its creation unless later lowered.
An eigenvariable $c$ created by a \scm{∀°} at depth $d$ has $\ell(c)=d+1$. The names that the unifier
itself opens when it goes under a binder have level $\infty$.
\end{definition}
These levels are Nadathur, Jayaraman and Kwon's tags minus one: their query starts at tag 1
(\cref{fig:rules-ui}), ours at depth 0.

The level is the linear trace of Miller's mixed prefix: a variable of level $d$ stands for an
existential quantifier that follows the first $d$ universal quantifiers of its branch. The intended
reading is Nadathur, Jayaraman and Kwon's~\cite{NadathurJayaramanKwon1995} (\cref{def:lkanren-ui}): a term may be substituted for $X$ only
if every constant in it has a level $\le \ell(X)$.

\begin{example}
In \scm{(fresh° (q) (∀° (x) (fresh° (y) G)))}, \scm{q} has level 0, the eigenvariable \scm{x} level 1,
and \scm{y} level 1. So \scm{y} may be bound to a term containing \scm{x}, while \scm{q} may not. This
is $\exists q\,\forall x\,\exists y.\,G$ read off the creation depths.
\end{example}

\begin{definition}[Order independence]\label{def:design-order}
A conjunction is \emph{order independent} when permuting its conjuncts changes at most the order of its
answers, not the set of answers (taken as data, with their residual constraints).
\end{definition}

Order independence is not an item of its own. It runs through all of them, because it is what
Nadathur, Jayaraman and Kwon prove for their tagged interpreter: the choice of the goal to reduce next
is irrelevant (their Prop.~4.1, \cref{sec:lkanren-impl}). A miniKanren conjunction is fair to no
conjunct, so any order-dependence is visible to the user. Most defects found by the reviews of
\cref{sec:design-reviews} were violations of it, \enquote{in one goal order of three}.

\subsection{The nine items}\label{sec:design-items}

\begin{definition}[A1: the scope of \texorpdfstring{$\forall$}{∀}]\label{def:design-a1}\leavevmode
\begin{enumerate}[label=(\alph*)]
\item \textbf{\rulename{generic}.} In a state of depth $d$, \scm{(∀° (x) G)} runs $G[c/x]$ at depth $d+1$, with
  $c$ a new eigenvariable of level $d+1$, and gives every answer state of $G$ back the depth $d$.
\item \textbf{Scope invariant.} In every state reached by the search, every binding $X:=t$ of the
  substitution satisfies the three binding conditions of \cref{def:lkanren-ui} (stated also in \cref{sec:hh-scope} and
  \cref{def:rules-binding}):
  (1)~$X$ does not occur in $t$; (2)~every eigenvariable occurring in $t$ has a level $\le \ell(X)$;
  (3)~every unbound variable $Y$ of $t$ has $\ell(Y)\le\ell(X)$ (\emph{lowering}; for a flexible
  subterm $Y\,\bar a$, the arguments not visible at $\ell(X)$ are pruned).
\item \textbf{Universality of the invariant.} (b) holds for every binding, whatever goal made it: the
  goals of the new module, micro's own \scm{=°}, every relation built on it (\scm{cons°}, the user's
  relations), and the unifier's internal bindings.
\end{enumerate}
\end{definition}

Intuitively, (a) is the book's \rulename{generic} rule \book{82} (\cref{def:hh-eigen}) and (b) is the reason why
it is sound: no eigenvariable can reach a variable created before its \scm{∀°}, directly or through an
alias, so a proof never assumes that a fresh name equals an older term. Restoring the depth on each
answer matters: a variable created \emph{after} the \scm{∀°} has exited must not see $c$, which a creation
timestamp alone would not ensure.

The invariant (b) is stated on single bindings, but it controls whole values. Suppose $X:=t$ and $Y$
occurs unbound in $t$; by (3), $\ell(Y)\le\ell(X)$, so if $Y$ is bound later to some $t'$, condition
(2) for $Y:=t'$ admits only eigenvariables of level $\le\ell(Y)\le\ell(X)$. By induction on the
chain of bindings, the fully walked value of $X$ never holds an eigenvariable of level $>\ell(X)$:
checking each binding once, when it is made, is enough. Condition (3) is what makes the induction go
through; without lowering, an alias created before the eigenvariable could receive it later.

Item (c) is what separates a sound design from a layered one, as
\cref{sec:design-d1} shows. The book's two classical queries \book{85} and two escapes through aliases:
\begin{scheme}
(show "∃Y∀x. Y = x  and  ∀x∃Y. Y = x"
      (list (°->list #f (fresh° (q) (∀° (x) (=° q x))))
            (°->list #f (fresh° (q) (∀° (x) (fresh° (y) (=° y x))) (=° q 'ok)))))

(show "escape through an alias (q = (y), y = x), and through micro's cons°: both fail"
      (list (°->list #f (fresh° (q) (∀° (x) (fresh° (y) (=° q (list y)) (=° y x)))))
            (°->list #f (fresh° (q) (∀° (x) (cons° x '() q))))))
\end{scheme}
\begin{answers}
;; ∃Y∀x. Y = x  and  ∀x∃Y. Y = x
(() ((λ () 'ok)))

;; escape through an alias (q = (y), y = x), and through micro's cons°: both fail
(() ())
\end{answers}
In the third query the binding \scm{y := x} is legal (\scm{y} has level 1), but \scm{q := (y)} was made
earlier with \scm{y} unbound; condition (3) lowered \scm{y} to level 0 at that moment (by binding it
to a fresh variable of level 0), so the later
\scm{y := x} violates (2). The fourth query never mentions the new module's unifier: \scm{cons°} is
micro's, and its \scm{=°} is checked all the same. The suite adds Nadathur, Jayaraman and Kwon's own
example (\cref{sec:hh-scope}, and the tutorial after \cref{def:lkanren-ui}), where lowering reaches a variable created by
backchaining on an assumed clause: $\exists x\forall y(q\,y\imp p\,x)$ with the clause
$\forall z(q\,z\imp p\,(d\,z))$ fails, as it must.

\begin{definition}[A2: \rulename{augment} and \rulename{backchain}]\label{def:design-a2}\leavevmode
\begin{enumerate}[label=(\alph*)]
\item \textbf{\rulename{augment}.} $\Sig;\Prog\longrightarrow D\imp G$ reduces to $\Sig;\Prog\cup\{D\}\longrightarrow
  G$: \scm{(⊃° (D ...) G)} runs $G$ with the clauses $D$ added to a context $\Gamma$ kept in the state, and
  gives every answer state of $G$ back the old $\Gamma$. Hypotheses are visible exactly during $G$, and
  the branches of an \scm{or°} do not see each other's hypotheses.
\item \textbf{Clause order.} $\Gamma$ is a list, extended at the front. A dynamic relation $p$ called on
  $\bar t$ is the disjunction of its clauses in $\Gamma$, most recent first, followed by its static body
  \book{78}.
\item \textbf{Variables of a clause.} The logic variables free in $D$ belong to the context: they are
  shared by every use of the clause and with the goal around it \book{86--88}. The clause's own
  universally quantified variables, \scm{(∀ (z) ...)}, are renamed at each use, at the current level
  (Nadathur, Jayaraman and Kwon's \enquote{new tagged instance}).
\item \textbf{Static relations are closed.} Assuming a clause for a relation not declared with
  \scm{define-dynamic-relation} is an error, not a silent no-op.
\end{enumerate}
\end{definition}

Point (b) is a statement about \emph{order}, which a miniKanren usually leaves to interleaving. The book
fixes it with two queries against the program \lp{p 1} \book{78}:
\begin{lprolog}
?- p 2 => p 3 => p X.
X = 3;
X = 2;
X = 1
?- (p 2 & p 3) => p X.
X = 2;
X = 3;
X = 1
\end{lprolog}
\begin{scheme}
(define-dynamic-relation (p° x) (=° x 1))
(show "clause order under augmentation (p.94)"
      (list (°->list/datum (fresh° (q) (⊃° ((p° 2)) (⊃° ((p° 3)) (p° q)))))
            (°->list/datum (fresh° (q) (⊃° ((p° 2) (p° 3)) (p° q))))))
\end{scheme}
\begin{answers}
;; clause order under augmentation (p.94)
((3 2 1) (2 3 1))
\end{answers}
(The label's \enquote{p.94} is the PDF page, \book{78}; see \cref{sec:intro-conventions}.) Each clause answers without
suspending, so the flat disjunction returns its disjuncts in order, and the order is the book's. Point
(c) is the subtle one: in \lp{pi z\ (q z => p (d z))} the variable \lp{z} is the clause's own, while a
free \lp{Y} in an assumed \lp{p Y} is the goal's. \Cref{sec:hh-kanren-context} explains why $\Gamma$ must
live in the state rather than in a Scheme parameter.

\begin{definition}[A3: \texorpdfstring{$\forall$}{∀} and \texorpdfstring{$\imp$}{⊃} together]\label{def:design-a3}
The book's type inference for untyped λ-terms (\cref{sec:lambdatree-typeof}), whose abstraction case is
\lp{typeof (abs M) (arr A B) :- pi x\ typeof x A => typeof (M x) B}, computes principal types: for a
closed term it answers the book's type up to the names of type variables, and it fails on terms that
have no simple type.
\end{definition}

This item exercises all of A1, A2 and A4 in one clause: \scm{∀°} makes the bound variable a scoped
constant, \scm{⊃°} assumes its type, \scm{(@ R x)} β-reduces the object-level body, and the type
variables \scm{A}, \scm{B}, created outside the \scm{∀°}, must never be bound to a type containing
\scm{x}. In the module's syntax:
\begin{scheme}
(define-dynamic-relation (typeof° m t)
  (cond°
    ((fresh° (M N B) (=° m `(app ,M ,N)) (typeof° M `(arr ,B ,t)) (typeof° N B)))
    ((fresh° (R A B) (=° m `(abs ,R)) (=° t `(arr ,A ,B))
       (∀° (x) (⊃° ((typeof° x A)) (typeof° (@ R x) B)))))))
\end{scheme}
Run on the combinators S and I and on $\lambda x.\,x\,x$, written in the object syntax with
\scm{ƛ} for the meta-level binder:
\begin{scheme}
(show "typeof S, I and λx.x x"
      (list (°->list/datum (fresh° (q) (typeof° `(abs ,(ƛ (x) `(abs ,(ƛ (y) `(abs ,(ƛ (z)
                                          `(app (app ,x ,z) (app ,y ,z)))))))) q)))
            (°->list/datum (fresh° (q) (typeof° `(abs ,(ƛ (x) x)) q)))
            (°->list/datum (fresh° (q) (typeof° `(abs ,(ƛ (x) `(app ,x ,x))) q)))))
\end{scheme}
\begin{answers}
;; typeof S, I and λx.x x
(((arr (arr α (arr β γ)) (arr (arr α β) (arr α γ)))) ((arr α α)) ())
\end{answers}
The third query fails by the occurs check: the type of \scm{x} would have to be an arrow type
containing itself.

\begin{definition}[A4: λ-tree syntax and β]\label{def:design-a4}\leavevmode
\begin{enumerate}[label=(\alph*)]
\item \textbf{α.} \scm{(ƛ (x) t)} builds an abstraction whose bound variable is a Scheme variable;
  terms built alike are \scm{equal?}, and unification is modulo α.
\item \textbf{β.} \scm{(@ (ƛ (x) t) s)} is $t[s/x]$ at construction; an application whose head becomes an
  abstraction through a later binding is reduced on the fly, in unification and in every deep walk.
\item \textbf{η.} Terms are kept η-short; η-expansion happens on demand when one side of an equation is
  an abstraction.
\item \textbf{Object-level substitution is meta-level β} (\cref{sec:lambdatree-beta}): the book's
  evaluators, \texttt{copy} and \texttt{subst} give the book's results.
\end{enumerate}
\end{definition}

\begin{scheme}
(show "α-equivalence: λa.a = λb.b, and the β-redex (λx.g x x) a"
      (list (°->list #f (fresh° (q) (=° (ƛ (a) a) (ƛ (b) b))))
            (°->list/datum (fresh° (q f) (=° f (ƛ (x) `(g ,x ,x))) (=° q (@ f 'a))))))
\end{scheme}
\begin{answers}
;; α-equivalence: λa.a = λb.b, and the β-redex (λx.g x x) a
(((λ (α) α)) ((g a a)))
\end{answers}
The first unification succeeds and binds nothing, so \scm{q} stays fresh. In the second, \scm{f} is
bound first and the application \scm{(@ f 'a)} is β-reduced when the deep walk of the answer meets it.
With call-by-name evaluation, the book's $(\lambda x\lambda w.w)\,\Omega$ \book{183} evaluates to
\lp{abs w\w} without ever evaluating Ω, because the substitution of the argument is a meta-level β that
drops it (\cref{ch:examples}).

\begin{definition}[A5: \texorpdfstring{$\Llam$}{Lλ} unification and delay]\label{def:design-a5}
Call \emph{pattern} for a variable $F$ an argument list $\bar a$ of $F\,\bar a$ whose members
head-normalise to distinct eigenvariables of level $>\ell(F)$ (\cref{def:hounif-llambda}, with levels in
place of the prefix).
\begin{enumerate}[label=(\alph*)]
\item \textbf{Patterns.} For an equation all of whose flexible subterms are patterns, unification
  terminates; it fails if and only if the equation has no unifier, and otherwise it binds a most general
  unifier (\cref{thm:hounif-mgu}). On the problems 8.G--8.K it gives the book's answers and mgus
  \book{222--225}. (The book does not number these problems; the labels 8.A--8.M are those of
  \cref{sec:hounif-traces}, in their order of appearance, as the test suite uses them; 8.M is the problem
  of \book{208}.)
\item \textbf{Dynamic $\Llam$} \book{225--226}. An equation outside $\Llam$ is neither failed nor
  enumerated: it is suspended, retried after every later binding, solved as soon as it becomes a
  pattern, and otherwise reported in the answer as a residual \scm{(assert (≐ u v))}, closed over the
  eigenvariables whose scope it outlived.
\end{enumerate}
\end{definition}

The level condition is Miller's \enquote{universally quantified inside the scope of $F$}. To see it
at work, compare $\exists F\,\forall x.\,F\,x=x$ with $\forall x\,\exists F.\,F\,x=x$. In the first,
$\ell(F)=0<1=\ell(x)$, so \scm{(@ F x)} is a pattern, and the only unifier is $F:=\abs{x}{x}$. In the second, $F$ is created
under the \scm{∀°}, so $\ell(F)=1=\ell(x)$ and \scm{(@ F x)} is not a pattern: $F$ may mention
$x$ as a constant, and both $\abs{z}{z}$ and $\abs{z}{x}$ solve the equation, so no unifier is most
general. The proposal solves the first; the second it suspends, accepting either solution when a
later goal supplies it:
\begin{scheme}
(pretty-print
 (list (°->list/datum (fresh° (q) (∀° (x) (=° (@ q x) x))))
       (°->list/datum (fresh° (q) (∀° (x) (fresh° (F) (=° (@ F x) x)
                                                    (=° F (ƛ (z) z)) (=° q 'id)))))
       (°->list/datum (fresh° (q) (∀° (x) (fresh° (F) (=° (@ F x) x)
                                                    (=° F (ƛ (z) x)) (=° q 'const)))))))
(pretty-print
 (°->pretty (fresh° (q) (∀° (x) (fresh° (F) (=° (@ F x) x) (=° q 'ok))))))
\end{scheme}
\begin{answers}
(((ƛ (x) x)) (id) (const))
((ok : (= (ƛ (x) (@ α x x)) (ƛ (x) x))))
\end{answers}
The last residual is instructive. When the \scm{∀°} exits, $F$ (of level 1) is \emph{raised},
$F:=\alpha\,x$ with $\ell(\alpha)=0$, and the equation is closed over $x$: $\forall x\,\exists F.\,F\,x=x$
has become $\exists\alpha\,\forall x.\,\alpha\,x\,x=x$, which is problem 8.M below.

Part (a) is Miller's theorem~\cite{Miller1991}, whose algorithm (\cref{def:hounif-elim}) prunes and raises; part (b) is
what Teyjus~2 does with the equations it cannot solve, after Dowek, Hardin, Kirchner and
Pfenning~\cite{DowekHardinKirchnerPfenning1996} (\cref{sec:hounif-dynamic}). The book's main
example, problem 8.H \book{223--224} (\cref{ex:hounif-8h}), $\forall w\,\forall x\,\forall y.\ f\,(U\,x\,y)=f\,(g\,(V\,y\,w))$
with $U$, $V$ quantified outside, is the λProlog query
\lp{pi w\ pi x\ pi y\ f (U x y) = f (g (V y w))}; it has the mgu
$U=\abs{x}{\abs{y}{g\,(V'\,y)}}$, $V=\abs{y}{\abs{w}{V'\,y}}$ (\cref{ex:hounif-8h}):
\begin{scheme}
(show "Example 8.H, pruning: (U V)" (°->list/datum (fresh° (q u v) (=° q (list u v)) (∀° (w x y) (=° `(f ,(@ u x y)) `(f (g ,(@ v y w))))))))
\end{scheme}
\begin{answers}
;; Example 8.H, pruning: (U V)
(((ƛ (x y) (g (@ α y))) (ƛ (y w) (@ α y))))
\end{answers}
The fresh \scm{α} is the book's $V'$. A non-pattern such as problem 8.M, \lp{pi x\ F x x = x}, has two
incomparable unifiers, so it waits; it is reported, and a later binding wakes it. Problem 8.A,
$F\,a=g\,a\,a$ with $a$ a constant, is not a pattern either (its argument is no eigenvariable):
\begin{scheme}
(show "non-patterns: 8.M and 8.A, suspended; 8.M woken by F := λa.λb.b"
      (list (°->pretty (fresh° (q) (∀° (x) (=° (@ q x x) x))))
            (°->pretty (fresh° (q) (=° (@ q 'a) '(g a a))))
            (°->list/datum (fresh° (q) (∀° (x) (=° (@ q x x) x))
                                       (=° q (ƛ (a b) b))))))
\end{scheme}
\begin{answers}
;; non-patterns: 8.M and 8.A, suspended; 8.M woken by F := λa.λb.b
(((α : (= (ƛ (x) (@ α x x)) (ƛ (x) x))))
 ((α : (= (@ α a) (g a a))))
 ((ƛ (a b) b)))
\end{answers}
(The first two answers, 8.M and 8.A, are printed by \scm{°->pretty} of the show module, the third
by \scm{°->list/datum}; the residual of 8.M is closed over \scm{x}, which has left its scope.)

\begin{definition}[A6: λProlog's semantics where Ma and Friedman differ]\label{def:design-a6}
On the equations by which Ma and Friedman~\cite{MaFriedman2021} motivate static closures (\cref{sec:lkanren-mf21}), the
answer is λProlog's:
(1)~$\abs{a}{X}=\abs{b}{b}$ and their conjunction containing it fail, since $\exists X\forall a.\,X=a$ is
false (\cref{ex:hounif-8j}); the paper answers $X=\mathtt{d0}$;
(2)~$\abs{a}{\abs{a}{X}}=\abs{c}{\abs{d}{c}}$ fails, with no false \enquote{true = false}, because shadowing
is lexical;
(3)~their Example~3, $\abs{a}{\abs{b}{X}}=\abs{c}{\abs{d}{Y}}$, gives $X=Y$;
(4)~an assumption is tried before the program: with a dynamic \scm{eq°}, their \texttt{assume} query gives
\scm{(orange apple)} where the paper prints \scm{(apple orange)};
(5)~their E12, $\forall A.\,A=F\,A$, gives $F=\abs{A}{A}$, as in the paper.
\end{definition}

\begin{scheme}
(show "λa.X = λb.b fails (the paper: X = d0); λa.λa.X = λc.λd.c fails; ∀A. A = F A gives F = λA.A"
      (list (°->list #f (fresh° (q) (=° (ƛ (a) q) (ƛ (b) b))))
            (°->list #f (fresh° (q) (=° (ƛ (a) (ƛ (a) q))
                                        (ƛ (c) (ƛ (d) c)))))
            (°->list/datum (fresh° (q) (∀° (a) (=° a (@ q a)))))))
\end{scheme}
\begin{answers}
;; λa.X = λb.b fails (the paper: X = d0); λa.λa.X = λc.λd.c fails; ∀A. A = F A gives F = λA.A
(() () ((ƛ (a) a)))
\end{answers}
\Cref{sec:lkanren-counter} discusses these equations; the point here is only that the proposal answers
them the way λProlog's quantifier reading dictates.

\begin{definition}[A7: the book's programs]\label{def:design-a7}
The book's programs over λ-tree syntax give the book's answers, in the book's order where the
interleaving search allows it: de Bruijn translation in both directions \book{193}; prenex normal forms,
the unique one and the five of a non-unique case \book{196--197}; \texttt{mappred} (Figs.~5.3--5.4);
palindromes over functional difference lists, the book's first nine answers \book{142}; Fischer's CPS
transformation, exactly the book's term \book{257}.
\end{definition}

These programs are the integration tests: each needs several of A1--A5 at once, and each has an answer
printed in the book to compare with. \Cref{ch:examples} walks through them.

\begin{definition}[A8: micro's stores on λ-terms]\label{def:design-a8}\leavevmode
\begin{enumerate}[label=(\alph*)]
\item \scm{≠°} decides modulo αβη and respects scope; its outcome does not depend on whether a binding
  comes before or after it. Under \scm{∀°} it reads an eigenvariable as \rulename{generic}'s new name by default
  (\lp{pi x\ not (q = x)} succeeds), and as any term of its universe when the parameter
  \scm{μkanren-λ-reading} is \scm{'extensional} (the reading of Moiseenko and of Jin et
  al.~\cite{Moiseenko2019,JinEtAl2021}).
\item \scm{absent°} looks for its datum in the parts of a λ-term, never in binder names, de Bruijn
  indices or the internals of an eigenvariable; on a flexible term it waits.
\item \scm{symbol°}, \scm{number°} and the user's types wait on a flexible term, and decide a rigid
  λ-term or an eigenvariable by the type's predicate.
\item No eigenvariable is ever printed free in an answer; residual constraints print as \scm{≐} and
  \scm{(∀ (x) ...)}; binders print renamed apart; ground answers evaluate back to live λ-terms.
\end{enumerate}
\end{definition}

λProlog has no disequality, \emph{absento} or types, so this item is not fidelity but coherence: micro's
stores must keep meaning what they mean once terms have binders. For instance:
\begin{scheme}
(show "≠ modulo αβη, flex-flex with an existential, a flexible absento, a type on a flexible term"
      (list (°->pretty (fresh° (q) (≠° q (ƛ (x) x))))
            (°->pretty (fresh° (q f g) (=° q (list f g))
                         (≠° (ƛ (x y) (@ f x)) (ƛ (x y) (@ g y)))))
            (°->pretty (fresh° (q F) (=° q F) (absent° 'c (@ F 'c))))
            (°->list/datum (fresh° (q F) (=° q F) (absent° 'c (@ F 'c))
                             (=° F (ƛ (z) 'b))))
            (°->list/datum (fresh° (q F) (=° q (@ F 'a)) (symbol° q)
                             (=° F (ƛ (z) z))))))
\end{scheme}
\begin{answers}
;; ≠ modulo αβη, flex-flex with an existential, a flexible absento, a type on a flexible term
(((α : (≠ α (ƛ (x) x))))
 (((α β) : (∀ (_3) (≠ (α β) ((ƛ (x) _3) (ƛ (y) _3))))))
 ((α : (absento c (@ α c))))
 ((ƛ (z) b))
 (a))
\end{answers}
Line by line. (1) $q\ne\abs{x}{x}$ is kept as a constraint, as micro keeps $q\ne 5$. (2) The trial
unification of $\abs{x}{\abs{y}{f\,x}}$ with $\abs{x}{\abs{y}{g\,y}}$ prunes both $f$ and $g$ to a
common constant function: the line says $\neg\exists H.\,f=\abs{x}{H}\wedge g=\abs{y}{H}$, printed in
the equivalent form $\forall H.\,(f\;g)\ne(\abs{x}{H}\;\abs{y}{H})$. The variable $H$ (\scm{_3}), which
the trial created, is quantified \emph{inside} the negation, which micro's prefixes cannot express.
(3) Whether \scm{c} is absent from $F\,\cst c$ depends on $F$, so the constraint waits on the whole
flexible term. (4) Once $F:=\abs{z}{\cst b}$, the argument is dropped and the constraint holds.
(5) \scm{(symbol° q)} waits on $F\,\cst a$ and holds once $F:=\abs{z}{z}$ makes it \scm{a}.
\Cref{ch:module} explains each of these answers in detail.

\begin{definition}[A9: conservativity]\label{def:design-a9}\leavevmode
\begin{enumerate}[label=(\alph*)]
\item micro's six suites (\texttt{microkanren}, \texttt{-aggregation}, \texttt{-untagged},
  \texttt{-show}, \texttt{-arith}, \texttt{-date}; 125 cases) pass with unchanged expectations, answer
  order included, against the edited micro, both with and without \scm{(aux kanren lambda)} loaded.
\item A first-order program pays at most a few per cent for the edit (measured in
  \cref{sec:design-cost}).
\end{enumerate}
\end{definition}

\Cref{tab:design-items} summarises the nine items and the evidence for each.

\begin{table}[ht]
\centering\footnotesize
\caption{The acceptance items and the hand-written cases that pin them (83 cases, 320 assertions in
\texttt{microkanren-lambda.scm}; 6 more cases in \texttt{microkanren-lambda-limits.scm} pin the known
limits).}
\label{tab:design-items}
\renewcommand{\arraystretch}{1.2}
\begin{tabularx}{\textwidth}{@{}l>{\raggedright\arraybackslash}X>{\raggedright\arraybackslash}p{2.3cm}r@{}}
\toprule
item & requirement & book & cases / assertions\\
\midrule
A1 & scope of ∀: \rulename{generic}, NJK's binding conditions on every binding & \book{82, 85} & 9 / 22\\
A2 & ⊃: \rulename{augment}, clause order, shared and renamed variables, closed static relations & \book{77--78, 86--88} & 8 / 21\\
A3 & ∀ and ⊃ together: type inference & \book{191} & 3 / 8\\
A4 & λ-tree syntax: α, β, η; object substitution by meta β & \book{177--199} & 4 / 9\\
A5 & $\Llam$ unification with mgus; dynamic $\Llam$ (delay) & \book{220--226} & 10 / 40\\
A6 & λProlog's answers where Ma--Friedman differ & --- & 5 / 13\\
A7 & the book's programs, the book's answers & \book{142, 193--197, 257} & 5 / 12\\
A8 & ≠, absento, types and reification on λ-terms & --- & 11 / 65\\
A9 & micro's suites unchanged; first-order cost & --- & 1 / 2, + 125 cases ×2\\
\midrule
R, R2, R3 & regressions of the three review rounds (\cref{sec:design-reviews}) & & 27 / 128\\
\bottomrule
\end{tabularx}
\end{table}

% ================================================================================================
\section{The design space}\label{sec:design-space}

\Cref{sec:lkanren-compare} put four candidate meanings of \scm{=°} on λ-terms side by side, without
choosing (\cref{tab:lkanren-options}). In one paragraph each:

\begin{description}
\item[A.] \textbf{Pattern unification with de Bruijn terms and levels}, after Miller, Nipkow, and
  Nadathur and Linnell~\cite{Miller1991,Nipkow1993,NadathurLinnell2005}; what Teyjus~2 and
  ELPI~\cite{DunchevEtAl2015} run. Decidable and unitary on $\Llam$; non-patterns are delayed.
\item[B.] \textbf{Nominal unification} (Urban--Pitts--Gabbay~\cite{UrbanPittsGabbay2004}; αKanren~\cite{ByrdFriedman2007}). Decidable, unitary with a freshness
  context. A different logic: names as data, no function variables, no β; noms are distinct constants,
  not eigenvariables, so \rulename{generic} still needs a separate mechanism; αKanren has no disequality.
\item[C.] \textbf{Static closures} (Ma--Friedman~\cite{MaFriedman2021}). Grafting semantics: $\abs{a}{X}=\abs{b}{b}$ succeeds,
  unlike λProlog; a workshop preview with no published algorithm.
\item[D.] \textbf{Huet's pre-unification, enumerated by the stream} (Huet~\cite{Huet1975}; Teyjus~1). Exact for full higher-order
  unification, but semi-decidable, with possibly infinitely many pre-unifiers.
\end{description}

\noindent The proposal retains A, for three reasons.

\begin{enumerate}
\item \textbf{λProlog's semantics.} The requirements are the book's. Option A is what the book's own
  implementations run, and with delayed non-patterns its behaviour on the book's examples is the book's
  (A5). Option C gives different answers on A6's equations by design; option B answers a different
  question.
\item \textbf{Decidability and most general unifiers.} Miller's theorem (\cref{thm:hounif-mgu}) makes
  unification on patterns terminating and unitary. This matches the shape of micro's \scm{=°}, which
  maps a state to \emph{at most one} state (\scm{✓°} or \scm{✗°}). Option D needs branching inside
  unification, so it cannot be an \scm{=°}; it can be added later as an explicit search over the
  suspended equations, which are its natural agenda (\cref{ch:roadmap}).
\item \textbf{Universe levels make \rulename{generic} cheap and exact.} The tagged interpreter of Nadathur,
  Jayaraman and Kwon is sound and complete, and needs one integer per variable and one comparison per
  binding (\cref{def:lkanren-ui}). The alternatives for ∀ alone are weaker or costlier: the explicit
  prefix of λKanren 2020~\cite{MaLuFriedman2020} is walked at each check and only grows (\cref{sec:lkanren-mk20}); the
  eigen-absento of Hemann, Byrd and Friedman~\cite{HemannByrdFriedman2013} costs $O(|B|)$ per check and is first-order only
  (\cref{sec:lkanren-eigen}); the end-of-scope check of \cref{sec:micro-forall} prunes late and lets
  constraints outlive their scope.
\end{enumerate}

Three secondary choices follow from A and are shared by every prototype below. Terms are
\emph{locally nameless}~\cite{McBrideMcKinna2004,Chargueraud2012} (de Bruijn indices inside bodies,
eigenvariables for free names), so α costs
nothing and terms built alike are \scm{equal?}. Raising is done lazily through levels rather than
eagerly as in the book's §4.4.1 \book{107}; the two differ in the shape of the mgu, not in its meaning
(\cref{ex:hounif-8k}). And non-patterns are delayed rather than rejected: rejecting them, as ELPI does
by default, would lose \texttt{subst R N (R N)} with \texttt{R} unknown, and the palindrome program.

What A leaves open is \emph{where} the code lives with respect to \micro, and that question was
answered by building it three times.

% ================================================================================================
\section{Three prototypes}\label{sec:design-prototypes}

Three complete prototypes were built against A1--A9. They share the representation (locally nameless
abstractions, applications, eigenvariables with levels) and the algorithmic core (variable elimination
with pruning and lowering, a store of suspended equations). They differ in architecture.
\Cref{tab:design-prototypes} compares them at a glance; the subsections describe each one, its strengths,
and the concrete defects the three reviews found.

\begin{table}[ht]
\centering\footnotesize
\caption{The three prototypes.}
\label{tab:design-prototypes}
\renewcommand{\arraystretch}{1.2}
\begin{tabularx}{\textwidth}{@{}p{2.3cm}>{\raggedright\arraybackslash}X>{\raggedright\arraybackslash}X>{\raggedright\arraybackslash}X@{}}
\toprule
& \textbf{d1, layered} & \textbf{d2, hooks} & \textbf{d3, native}\\
\midrule
edit to micro & none (byte-identical binary) & +105/−25 lines: a field and seven hooks & micro grows from 713 to 1217 lines; a 931-line diff\\
λ state & a record bound to the reserved variable 0 & a seventh field \texttt{X} & \texttt{X} holds U, E, P; levels on the variable record\\
unifier & a second one, \texttt{=λ°}, beside micro's & micro's, plus an \texttt{unify} hook & micro's, rewritten\\
runner & its own, \texttt{λrun} & micro's \texttt{°->list} & micro's \texttt{°->list}\\
suite & 29 cases, 137 assertions, hand-written & 42 cases, 102 assertions, hand-written & 14 cases, 116 assertions, a generated snapshot\\
rank & last for all three reviews & first for two reviews & first for the pedagogy review\\
\bottomrule
\end{tabularx}
\end{table}

\subsection{d1: a layer with no edit to micro}\label{sec:design-d1}

\paragraph{Architecture.} d1 takes \cref{sec:micro-noedit} at its word: it imports micro unchanged and
adds a pure library. Its whole context (current level, level boundaries and overrides, the clause
database, delayed equations, pending constraints) is a record bound to the reserved variable 0, the
\enquote{S-register} of Trick~2, overwritten by a persistent \scm{update/sbral} in $O(\log|S|)$, so
each branch of an \scm{or°} keeps its own context; the query variable is variable 1. Since micro's binding site is closed, d1 brings its own unifier \scm{=λ°}, its own
constraints \scm{≠λ°} and \scm{absentλ°}, and its own runner \scm{λrun}. Levels need no slot in the
variables: micro allocates indices monotonically, so a variable's level at creation is a step function
of its index, recorded as boundaries pushed by \scm{∀°}; lowering is recorded as overrides.

\paragraph{Strengths.} Micro's behaviour and speed are untouched by construction. d1's \scm{absent°}
was the most precise on flexible terms, it offered the extensional reading of ≠ under ∀, it had
an αKanren-style printer, a data runner comparable with \scm{equal?}, and a table of every place where
purity leaks, each pinned by a test. All these were grafted into the final proposal.

\paragraph{Defects.} The leaks are the architecture's, not the code's. Micro's \scm{=°}, and every
relation built on it, binds without d1 seeing the binding, so A1(c) fails:
\begin{scheme}
; d1-layered, its own runner
(pretty-print (λrun (q -1) (∀° (x) (cons° x '() q))))
(pretty-print (λrun (q -1) (∀° (x) (=λ° q x))))
\end{scheme}
\begin{answers}
((x₂))
()
\end{answers}
With d1's own unifier the escape is caught; with micro's \scm{cons°} the eigenvariable \scm{x₂} leaks
into the answer. The proposal answers \scm{()} to the first query (A1, above). The reviews found four
more defects: the module's goals do not run under \scm{°->list} (they need the register);
conjunction is not commutative, because d1 read a list headed by a symbol as a rigid application in
spine notation, conflating lists with λ-tree applications, and together with η this made the answers of
some conjunctions depend on the order of the conjuncts; d1 duplicates a large part of micro; and \scm{=λ°} costs
4.4~times micro's \scm{=°} on search. Its clause database was a linear list, on which the reviews
measured slowdowns of 18~times. d1 was ranked last by all three reviews.

\subsection{d2: an extension protocol in micro}\label{sec:design-d2}

\paragraph{Architecture.} d2 makes the smallest principled edit that opens the closed points of
\cref{tab:micro-open}, and builds the whole λ layer as a client module on top of it. The state gets a
seventh field \texttt{X}, which micro copies and never reads. A record of procedures, registered
globally, has hooks for \texttt{term?}, \texttt{unify}, \texttt{post}, \texttt{diseq}, \texttt{value},
\texttt{repr} and \texttt{forms}; micro calls them at the matching points. The edit was +105/−25 lines
(counted on its \texttt{micro.diff}), the module 708 lines.

\paragraph{Strengths.} The reviews found no soundness defect in d2. It kept terms η-short as an
invariant, so both $\abs{x}{q\,x}=q$ and $F\,x=G\,x\Rightarrow F=G$ hold; it had a \enquote{soft}
occurs check that suspends $F\,x=G\,(F\,x)$ instead of failing it (it holds with $G:=\abs{z}{z}$); its
suite was hand-written with book page references; its printer was the only one that renamed binders
correctly; and the edit cost at most 2.8\% on micro's benchmarks.

\paragraph{Defects.} The reviews found seven, listed in \cref{tab:design-d2}. None is a soundness
defect; most are imprecisions of the stores on λ-terms and one is a cost.

\begin{table}[ht]
\centering\footnotesize
\caption{The defects of d2 found by the reviews, and their fixes (design notes, §10.2).}
\label{tab:design-d2}
\renewcommand{\arraystretch}{1.2}
\begin{tabularx}{\textwidth}{@{}l>{\raggedright\arraybackslash}X>{\raggedright\arraybackslash}X@{}}
\toprule
& defect & fix in the proposal\\
\midrule
D1 & a suspended disequality printed as \texttt{(deny (and))} & the denied equation, \texttt{(deny (≐ u v))}\\
D2 & \texttt{absent°} over-strict on flexible terms & the \texttt{parts} hook: a flexible term is watched whole\\
D3 & trivially true residuals kept ($F\,x\,x=F\,x\,x$) & delay drops α-equal sides\\
D4 & Miller's n7 left suspended and printed with \texttt{\_2} & discharge at ∀ exit; the \texttt{terms} hook names residual variables\\
D5 & registering the extension slowed ≠°-heavy first-order code by 45--50\% & migration only when a new binding holds a λ-term\\
D6 & goal suffix on term constructors, levels hidden in property lists, λ-terms not \texttt{equal?} & \texttt{ƛ}, \texttt{@}; plain records thanks to \texttt{parts}, \texttt{value}, \texttt{atom?}; interned hints\\
D7 & \texttt{symbol°}/\texttt{number°} failed eagerly on flexible terms; \texttt{⊃°} on a static relation silently ignored & the \texttt{tag} hook; an error for non-dynamic relations\\
\bottomrule
\end{tabularx}
\end{table}

One more defect turned up while grafting: $F=\abs{x}{G\,(F\,x)}$, reached through η-short construction,
failed micro's occurs check before any hook could see it. That is why the final protocol has an
\texttt{occurs} hook.

\subsection{d3: a native engine}\label{sec:design-d3}

\paragraph{Architecture.} d3 integrates λ-terms into micro itself, in the style of ELPI and Teyjus. The
variable record gets a level slot; \scm{μkanren-unify} knows abstractions, eigenvariables and flexible
applications; the binding site does the occurs check, the scope check and lowering in one walk; \scm{∀°}
and \scm{⊃°} are core goals; the current universe, the delayed equations and the program are packed in
one extra state field. Micro grew from 713 to 1217 lines.

\paragraph{Strengths.} The pedagogy review ranked d3 first for three reasons: one vocabulary, readable
answers in \scm{ƛ}, \scm{@} and \scm{≐}, and answers usable as code. d3 also discharged residual
constraints at ∀ exit by raising and closing them, and it was the fastest: type inference of the Church
numeral 400 took 105~ms, against d2's 441~ms.

\paragraph{Defects.} The reviews found a crash in \scm{μkanren-verify-D+} on a disequality prefix whose
trial unification delays, and an η-incompleteness that made \scm{≠°} unsound: after
$\abs{x}{F\,x}\ne\abs{x}{G\,x}$, the binding $G:=F$ was accepted.
\begin{scheme}
(pretty-print (°->list #f (fresh° (q F G) (=° q (list F G)) (≠° (ƛ (x) (@ F x)) (ƛ (x) (@ G x))) (=° G F))))
\end{scheme}
Run against d3, then against the proposal:
\begin{answers}
((λ (α) (cons α (cons α '()))))
\end{answers}
\begin{answers}
()
\end{answers}
The answer of d3 is a pair $(F,F)$, which violates the disequality it was asked to keep. Besides, its
suite was a snapshot generated from the actual answers rather than written from the book, and its diff
carried trailing-whitespace churn.

\subsection{The verdict}\label{sec:design-verdict}

The three reviews gave aggregate scores of \textbf{23} to d2, \textbf{21} to d3 and \textbf{15.5} to
d1. Two reviews ranked d2 first, the pedagogy review ranked d3 first, and all three ranked d1 last. The
choice of d2 as the base rests on four facts.

\begin{itemize}
\item \textbf{Size and scope of the edit.} The final protocol is +235/−41 lines of micro (713 to 907
  lines; 214 of the added lines are not blank), comments included, all generic: nothing in micro mentions λ, levels or clauses. d3 special-cases every walker.
  The same protocol could carry another store, such as nominal terms or finite domains.
\item \textbf{Soundness} is equal. d3's eager binding rule and the λ layer's eager \texttt{post} check
  enforce the same binding conditions; every scope attack of the reviews fails correctly on both, and on
  the final version.
\item \textbf{Correctness} favours d2: d3 had a crash and an unsound \scm{≠°}, d2 neither.
\item \textbf{Pedagogy.} d3's three advantages were grafted onto d2 (one vocabulary of \scm{ƛ}, \scm{@},
  \scm{≐}; readable residuals; ground answers as live terms through the \texttt{environment} hook), so
  the pedagogical argument for the native design no longer holds.
\end{itemize}

The reviews also recommended about twenty grafts from d1 and d3. Each was either done or explicitly
declined; \cref{tab:design-grafts} lists the main ones.

\begin{table}[ht]
\centering\footnotesize
\caption{Grafts recommended by the reviews (a selection from the design notes, §10.3).}
\label{tab:design-grafts}
\renewcommand{\arraystretch}{1.2}
\begin{tabularx}{\textwidth}{@{}l>{\raggedright\arraybackslash}X>{\raggedright\arraybackslash}p{4.6cm}@{}}
\toprule
from & graft & status\\
\midrule
d3 & discharge residual constraints at ∀ exit by raising and closing & done\\
d3 & \texttt{≐} exported, residuals as \texttt{(assert (≐ …))}, merged binders \texttt{(ƛ (x y) …)} & done\\
d3 & a first-order fast path for the scope check & done (\texttt{μkanren-plain?}, \texttt{top}, the \texttt{atom?} hook)\\
d3 & ground answers as live terms & done (\texttt{environment} hook)\\
d3 & the 2--4× faster λ machinery & partly: Church 400 from 441 to 254 ms; d3: 105 ms\\
d3 & var--var bindings oriented by level & declined: a hook on micro's second-hottest path\\
d3 & levels on the variable record & declined: \texttt{freshª} would need the universe, so micro would become λ-aware\\
d1 & precise \texttt{absent°} on flexible terms & done\\
d1 & the extensional reading of ≠ under ∀, as an option & done (\texttt{μkanren-λ-reading})\\
d1 & a data runner, an αKanren-style view & done (\texttt{°->list/datum}, \texttt{°->pretty})\\
d1 & a limits table, each row pinned by a case & done\\
d1 & truly nameless λ-records, \texttt{equal?} modulo α & partly: interned hints; α-variants with different hints are not \texttt{equal?}\\
\bottomrule
\end{tabularx}
\end{table}

The two declined grafts share a reason. Both would put knowledge of levels into micro's core: the
var--var clause of the unifier, its hottest path after pair--pair, and the variable constructor
\scm{freshª}, which would need the current universe kept in the extension's data. The protocol's
discipline is that micro stays λ-agnostic; for the orientation, the lowering done by the scope check
yields the same mgu up to renaming.

% ================================================================================================
\section{The extension protocol}\label{sec:design-protocol}

The protocol is the part of the proposal that edits \micro. It is generic: it lets a client module own
a class of records (its \emph{terms}), keep per-branch data in the state, and be consulted at each point
where micro would otherwise treat those records as plain slot vectors. This section gives it in full.
Excerpts are verbatim from \texttt{micro.diff}; lines starting with \texttt{+} are added, with
\texttt{-} removed, and with a space unchanged context. Comment lines of the diff are left out of some
excerpts, as is the removed body of \scm{=°}.

\subsection{The seventh field}\label{sec:design-X}

The state gets a field \texttt{X}, the extensions' slot. Micro never reads it: it copies it at its two
positional constructions, the binding site \scm{μkanren-state-update} and \scm{freshª}, and
\scm{μkanren-state-with} accepts \texttt{X:}. The positional constructor keeps its six-argument
signature, with \texttt{X} optional, so the positional states of \texttt{test/microkanren.scm} still
work.
\begin{scheme}
-  (define-record μkanren-state vars-count S D A T tags)
+  (define-record-type μkanren-state
+    (make-μkanren-state/X vars-count S D A T tags X)
+    μkanren-state?
+    (vars-count μkanren-state-vars-count) (S μkanren-state-S) (D μkanren-state-D) (A μkanren-state-A)
+    (T μkanren-state-T) (tags μkanren-state-tags) (X μkanren-state-X))
+
+  (define (make-μkanren-state vars-count S D A T tags #!optional (X '())) (make-μkanren-state/X vars-count S D A T tags X))
\end{scheme}
and in \scm{freshª}:
\begin{scheme}
-             (s*  (make-μkanren-state (add1 vc) S D A T tags)))
+             (s*  (make-μkanren-state/X (add1 vc) S D A T tags (μkanren-state-X s))))
\end{scheme}
\texttt{X} is \scm{'()} until an extension writes it, which the λ layer reads as \enquote{everything
empty}; so \scm{°->§} keeps its fixed empty initial state. The λ layer stores in \texttt{X} one record
with the current depth, the highest depth reached, the variables' levels, the context $\Gamma$, the
suspended constraints E and the λ-disequalities N (\cref{ch:module}). Since S is persistent and
\texttt{X} is copied forward by every goal, the data is branch-local for free: each branch of an
\scm{or°} sees its own $\Gamma$, as A2(a) requires.

\subsection{Extensions and the registry}\label{sec:design-registry}

An extension is a record built by
\begin{scheme}
  (define (make-μkanren-extension name #!key (term? #f) (atom? #f) (unify #f) (occurs #f) (post #f) (diseq #f) (tag #f) (parts #f)
                                             (value #f) (repr #f) (terms #f) (forms #f) (environment #f))
\end{scheme}
and \scm{μkanren-extension-register!} adds it to a global list. Registering again under the same name
replaces the old extension, so reloading a module is harmless. Dispatch is by term: the extension that
owns a pair of terms is the first whose \texttt{term?} holds of either.
\begin{scheme}
+  (define (μkanren-extension/term u v)
+    (let loop ((E μkanren-extensions))
+      (cond
+        ((null? E) #f)
+        ((let1 (term? (μkanren-extension-term? (car E))) (and term? (or (term? u) (term? v)))) (car E))
+        (else (loop (cdr E))))))
\end{scheme}
The call sites are guarded by \scm{(pair? μkanren-extensions)} (those of \texttt{forms} and
\texttt{environment}, in reification and in the runner, just map over the empty list), so with no
extension registered a hook costs one \scm{pair?} test, and none sits on micro's hottest path (the variable clauses and the
pair--pair clause of the unifier). The registry is global rather than per state because \scm{°->list}
starts from \scm{μkanren-state-empty}: a per-state hook list would have to be carried by every initial
state.

The protocol's whole exported surface is small: \scm{μkanren-state-X} (and the \texttt{X:} keyword of
\scm{μkanren-state-with}), \scm{make-μkanren-extension}, \scm{μkanren-extension?},
\scm{μkanren-extension-name}, \scm{μkanren-extension-register!}, and the two primitives
\scm{μkanren-post-unify} and \scm{μkanren-state-forget} of \cref{sec:design-forget}. The hook
accessors and the dispatchers (\scm{μkanren-extension/term} and its siblings) stay internal to micro.

The λ layer registers itself once, at load time:
\begin{scheme}
  (define μkanren-λ-extension
    (make-μkanren-extension 'λProlog
      term?: μkanren-λ-term? atom?: (λ (u) (or (μkanren-eigen? u) (μkanren-bvar? u))) unify: μkanren-λ-unify occurs: μkanren-λ-occurs post: μkanren-λ-post diseq: μkanren-λ-diseq
      tag: μkanren-λ-tag parts: μkanren-λ-parts value: μkanren-λ-value repr: μkanren-λ-repr
      terms: μkanren-λ-terms forms: μkanren-λ-forms environment: μkanren-λ-environment))

  (μkanren-extension-register! μkanren-λ-extension)
\end{scheme}

\subsection{The hooks}\label{sec:design-hooks}

\Cref{tab:design-hooks} lists the thirteen hooks. Compared with d2's seven, the final protocol adds
\texttt{atom?}, \texttt{occurs}, \texttt{tag}, \texttt{parts}, \texttt{terms} and \texttt{environment},
each answering a defect or a graft of the reviews (\texttt{atom?} the first-order fast path,
\texttt{environment} live ground answers, the others the defects of \cref{tab:design-d2} and the
occurs failure found while grafting). The last column says which closed point of
\cref{tab:micro-open} each hook opens. Then a paragraph per hook gives its contract and its call site.
\newpage
\begin{longtable}{@{}>{\raggedright\arraybackslash}p{2.0cm}>{\raggedright\arraybackslash}p{5.5cm}>{\raggedright\arraybackslash}p{3.3cm}>{\raggedright\arraybackslash}p{2.5cm}@{}}
\caption{The hooks of an extension. $u,v$: walked terms; $s$: a state; \emph{as}: the bindings made so
far, as \texttt{(α . t)} pairs; $\to$ gives the result, \textbar\ separates alternatives.}\label{tab:design-hooks}\\
\toprule
hook & arguments $\to$ result; contract & called by & opens\\
\midrule\endfirsthead
\toprule
hook & arguments $\to$ result; contract & called by & opens\\
\midrule\endhead
\texttt{term?} & \texttt{(u)}: $u$ is one of the extension's records & dispatch & ---\\
\texttt{atom?} & \texttt{(u)}: $u$ holds no variable & micro's occurs check, which then does not look inside $u$ & the binding site\\
\texttt{unify} & \texttt{(u v s as)} $\to$ \texttt{(s* as*)} \textbar\ \texttt{\#f}: a pair with one extension term, neither a variable & \texttt{μkanren-state-unify}: the structure--structure clause and the final \texttt{else} & the unification \texttt{cond}\\
\texttt{occurs} & \texttt{(α v s as)} $\to$ \texttt{(s* as*)} \textbar\ \texttt{\#f}: micro's occurs check failed on $α:=v$ & \texttt{μkanren-state-update}, failure path only & the binding site\\
\texttt{post} & \texttt{(s* as)} $\to$ \texttt{(s** as**)} \textbar\ \texttt{\#f}: runs after the bindings of \texttt{=°}, before D/A/T are verified; may bind more & \texttt{μkanren-post-unify}, called by \texttt{=°} & the post-\texttt{=°} chain\\
\texttt{diseq} & \texttt{(u v s)} $\to$ \texttt{\#f} \textbar\ state \textbar\ \texttt{'fail}: may claim a disequality & \texttt{≠°} & \texttt{≠°} (store D)\\
\texttt{tag} & \texttt{(u tag s)} $\to$ state \textbar\ \texttt{\#f}: a type on its term & \texttt{symbol°}, \texttt{number°}, user types & the A store\\
\texttt{parts} & \texttt{(u s [projecting])} $\to$ \texttt{'stuck} \textbar\ list: where \texttt{absent°} looks into its term; when projecting, the parts whose constraints an answer keeps & \texttt{absent°}, the re-verification of T, \texttt{μkanren-project-T} & the T store; reification\\
\texttt{value}, \texttt{repr} & \texttt{(u walk s)}: the deep value; the expression that rebuilds it & \texttt{μkanren-state-find/value}, \texttt{/repr} & reification\\
\texttt{terms} & \texttt{(s)} $\to$ terms: the terms of its residual constraints & \texttt{μkanren-project} & reification\\
\texttt{forms} & \texttt{(s s*)} $\to$ forms: the forms its constraints add to an answer & \texttt{μkanren-project} & reification\\
\texttt{environment} & \texttt{()}: where \texttt{°->list/ground} evaluates answers & \texttt{°->list} & the runner\\
\bottomrule
\end{longtable}

\paragraph{\texttt{term?}} identifies the extension's records. Every dispatch by term goes through it
(\texttt{unify}, \texttt{tag}, \texttt{parts}, \texttt{value}, \texttt{repr}); the other hooks are
consulted on every registered extension, so an extension without \texttt{term?} could still use them.
For the λ layer, the four records are abstractions,
applications, eigenvariables and de Bruijn indices.

\paragraph{\texttt{unify}} is consulted when micro's unifier meets a pair of walked terms, one of which
is the extension's and neither of which is a variable. Two clauses of the unification \scm{cond} hand
such pairs over: the structure--structure clause (two records, or a record and a vector) and the final
\scm{else} (an extension term against a pair or an atom).
\begin{scheme}
         ((and (μkanren-structure? u*) (μkanren-structure? v*))
-          (μkanren-state-unify (μkanren-structure->list u*) (μkanren-structure->list v*) s assocs))
+          (cond
+            ((and (pair? μkanren-extensions) (μkanren-extension/term u* v*)) => (λ (e) ((μkanren-extension-unify e) u* v* s assocs)))
+            (else (μkanren-state-unify (μkanren-structure->list u*) (μkanren-structure->list v*) s assocs))))
         ; the hottest path: a unification answers `(s* assocs*)` or #f, no need for a checked match.
         ((and (pair? u*) (pair? v*))  (let1 (r (μkanren-state-unify (car u*) (car v*) s assocs))
                                         (and r (μkanren-state-unify (cdr u*) (cdr v*) (car r) (cadr r)))))
+        ; an extension term against a pair or an atom.
+        ((and (pair? μkanren-extensions) (μkanren-extension/term u* v*)) => (λ (e) ((μkanren-extension-unify e) u* v* s assocs)))
         (else #f))))
\end{scheme}
The variable clauses, above this excerpt, and the pair--pair clause are untouched. The hook may call
\scm{μkanren-state-unify} back on parts: the λ layer's ξ rule opens both bodies with one new name and
unifies them with micro's own unifier. So a list nested in a λ-term is still unified by micro, and a
λ-term nested in a list is still handed to the hook. Before the edit, two records unified slot by
slot, so an abstraction would have been compared by its hint and its raw body.

\paragraph{\texttt{atom?}} tells micro's occurs check that a record holds no variable, so the check need
not look into its slots. For the λ layer, eigenvariables and indices are atoms. Without it micro rebuilt
every eigenvariable as a slot list at each binding; adding it cut type inference of the Church numeral
400 from 311 to 254~ms (an 18\% gain on the λ workloads).
\begin{scheme}
-                ((μkanren-structure? w) (O (μkanren-structure->list w)))
+                ((μkanren-structure? w) (and (not (μkanren-extension/atom? w)) (O (μkanren-structure->list w))))
\end{scheme}

\paragraph{\texttt{occurs}} gets the last word when micro's occurs check fails on $α:=v$ and $v$ may hold
extension terms. Inside a λ-term an occurrence may be only apparent: $q = F\,q$ holds with
$F:=\abs{z}{z}$, and $F=\abs{x}{G\,(F\,x)}$ with $G:=\abs{z}{z}$. The binding site, closed in
\cref{tab:micro-open}, is opened on its failure path only:
\begin{scheme}
+    (cond
+      ((or unsafe (not (occur? v)))
+        (μkanren-state-match ((vc S D A T tags) s)
+          (let* ((S* (μkanren-update/sbral α v S))
+                 (s* (make-μkanren-state/X vc S* D A T tags (μkanren-state-X s)))
+                 (assocs* (cons `(,α . ,v) assocs)))
+            `(,s* ,assocs*))))
+      ; an occurrence inside an extension's term may be only apparent: the extension decides.
+      ((pair? μkanren-extensions) (μkanren-extension/occurs α v s assocs))
+      (else #f)))
\end{scheme}
The λ layer answers by β-normalising $v$ (if it is $α$ itself the equation is trivial) and otherwise by
variable elimination, which fails on a rigid occurrence and suspends one in the argument of a flexible
term (\cref{ch:module}).

\paragraph{\texttt{post}} is the hook that carries the scope check. It runs after the bindings that one
\scm{=°} made and \emph{before} D, A and T are verified again; it receives those bindings, may bind more
(lowering, pruning, waking suspended equations), and returns the extended bindings, so that micro
re-verifies the stores they touch. The post-\scm{=°} chain, closed and wired for D/A/T before the edit,
becomes an exported procedure:
\begin{scheme}
+  (define (μkanren-post-unify s* assocs)
+    (let loop ((E μkanren-extensions) (s* s*) (assocs assocs))
+      (cond
+        ((null? E) (cond
+                     ((μkanren-verify-D/touched (μkanren-state-D s*) s* assocs)
+                       => (μ D* (μkanren-verify-D/post D* (μkanren-state-A s*) (μkanren-state-T s*) s*)))
+                     (else #f)))
+        ((μkanren-extension-post (car E)) => (λ (post)
+                                                (match/first (post s* assocs)
+                                                  (#f #f)
+                                                  ((,s** ,assocs*) (loop (cdr E) s** assocs*)))))
+        (else (loop (cdr E) s* assocs)))))
\end{scheme}
and \scm{=°} calls it:
\begin{scheme}
   (define ((=° u v) s)
+    (define post°
+      (λ1-match/first
+        ((,s* ,assocs)  (cond
+                          ((μkanren-state-equal? s s*) (✓° s*))
+                          ((μkanren-post-unify s* assocs) => ✓°)
+                          (else (✗° s*))))))
+    (cond
+      ((μkanren-state-unify u v s '()) => post°)
+      (else (✗° s))))
\end{scheme}
Read \scm{post°} clause by clause. If raw unification returned the very state it was given
(\scm{μkanren-state-equal?} is the \scm{eq?} test of \cref{prop:micro-eq}), nothing was bound, so no
hook runs and no store is re-verified; this is why invariant (I1) of \cref{sec:design-invariants}
matters. Otherwise the bindings \scm{assocs} go through every extension's \texttt{post}, in
registration order, each one receiving the state and bindings the previous one answered; then D is
re-verified on the touched prefixes and A and T after it. A \scm{#f} anywhere fails the \scm{=°}.

This is how A1(c) is met without opening the binding site: \emph{every} binding reaches the store
through some \scm{=°}, micro's included, and every \scm{=°} ends in \scm{μkanren-post-unify}. The binding
is checked a moment after it is made, but before any other goal can see it. Running \texttt{post}
before the D/A/T verification, not after, matters for two reasons: its own bindings must reach
\scm{μkanren-verify-D/touched}, and the λ layer moves a disequality prefix that has come to hold a
λ-term out of D \emph{before} micro re-unifies it.

\begin{pitfall}[A goal that unifies on its own must end in \texttt{μkanren-post-unify}]
The guarantee \enquote{every binding is checked} holds because every goal of micro that commits a
binding to the search state is an \scm{=°}; micro's other calls of raw unification are trials whose
states are thrown away (in \scm{≠°}, in \scm{absent°}, in subsumption) or reification. A goal of an extension that calls the raw \scm{μkanren-state-unify} and
returns its state bypasses \texttt{post} exactly as Trick~1's \scm{bind/raw} bypasses D/A/T
(\cref{sec:micro-noedit}): an eigenvariable could reach an older variable unchecked. Such a goal must
pass the state and the exact new bindings to \scm{μkanren-post-unify}, which is why it is exported; the
λ layer does so at \scm{∀°} exit (\cref{sec:design-forget}).
\end{pitfall}

\paragraph{\texttt{diseq}} lets an extension claim a disequality; micro's original body becomes
\scm{μkanren-≠}.
\begin{scheme}
+  ; an extension may take the disequality over, when its terms are involved.
   (define ((≠° u v) s)
+    (cond
+      ((and (pair? μkanren-extensions) (μkanren-extension/diseq u v s)) => (λ (r) (if (eq? r 'fail) (✗° s) (✓° r))))
+      (else (μkanren-≠ u v s))))
+
+  (define (μkanren-≠ u v s)
\end{scheme}
\texttt{\#f} leaves $u\ne v$ to the store D, a state keeps it in the extension's store, \scm{'fail}
says that $u$ and $v$ are already equal. A separate hook is needed because a D prefix is a list of
bindings: the mgu of a pattern trial can contain variables the trial itself created, which are
existential inside the negation (the \scm{∃} of item A8), and a prefix cannot express that.

\paragraph{\texttt{tag}} decides a type goal (\scm{symbol°}, \scm{number°}, a type made by
\scm{define-μkanren-tag}) on an extension term. Before the edit such a goal failed on every record that
is not of the type, so \scm{(symbol° (@ F 'a))} failed even though $F:=\abs{z}{z}$ would make it a
symbol; now the λ layer makes it wait.
\begin{scheme}
+      ((,u* ⊣ (μkanren-extension/hook u* μkanren-extension-tag)) (cond
+                                                                  (((μkanren-extension/hook u* μkanren-extension-tag) u* tag s) => ✓°)
+                                                                  (else (✗° s))))
\end{scheme}

\paragraph{\texttt{parts}} says where \scm{absent°} looks into an extension term: a list of parts, or
\scm{'stuck} when the term cannot be looked into yet, in which case \scm{absent°} watches it whole, like a
variable. Micro walks record slots, so without the hook \scm{(absent° 0 q)} would reject λx.x because
of the index 0 in its body, as it did in d1. With \texttt{projecting} true, the same hook tells the projection of T on an answer which
parts' constraints the answer keeps, instead of splitting a λ-term into its raw slots.
\begin{scheme}
+  (define (μkanren-absent-parts u s)
+    (cond
+      ((μkanren-extension/hook u μkanren-extension-parts) => (λ (parts) (parts u s)))
+      (else (μkanren-compound->list u))))
\end{scheme}
Its call sites are \scm{absent°}, \scm{μkanren-verify-T+} (both read \scm{'stuck} as \enquote{watch}) and
\scm{μkanren-project-T}, through the twin \scm{μkanren-project-parts}.

\paragraph{\texttt{value} and \texttt{repr}} give the deep value of an extension term and the Scheme
expression that rebuilds it; \texttt{walk} is passed so that the hook can recurse on the parts. They are
called by \scm{μkanren-state-find/value} and \scm{/repr}, so β-normalisation applies wherever micro
deep-walks: answers, the stores in answers, \scm{project°} and aggregation.
\begin{scheme}
+            ((and (record-instance? w*) (not (μkanren-var? w*)))
+              (cond
+                ((and (pair? μkanren-extensions) (μkanren-extension/term w* w*)) => (λ (e) ((μkanren-extension-value e) w* A s)))
+                (else (apply make-record-instance (map A (μkanren-structure->list w*))))))
\end{scheme}

\paragraph{\texttt{terms} and \texttt{forms}} give the extension a say in the answer. \texttt{terms}
returns the terms of its residual constraints whose variables the answer should name: micro reifies them
after the answer, so a residual prints \scm{α}, \scm{β} instead of \scm{_2}. \texttt{forms} returns the
forms its constraints add to the answer's body, after micro's own.
\begin{scheme}
+    (let* ((w* (μkanren-state-find/value w s))
+           (extra (μkanren-extension/terms s)))
+      ; the variables of the extensions' residual constraints are named too, after those of the answer.
+      (match1/first ((,s* ,c ,vars-reversed ,assocs) (μkanren-state-reify (if (null? extra) w* (cons w* extra)) s))
\end{scheme}
\begin{scheme}
+               (X**   (μkanren-extension/forms s s*))
+               (body  (μkanren-form w** D** A** T** X** s*)))
\end{scheme}

\paragraph{\texttt{environment}} names the environment in which \scm{°->list/ground} evaluates answers.
The λ layer gives its own module environment, where \scm{ƛ}, \scm{@} and \scm{≐} are bound, so a ground
answer evaluates back to a live λ-term that can be applied again (the pedagogy graft from d3).
\begin{scheme}
-                         (let ((E (eval `(lambda ,(cadr expr) ,(last expr)) (μkanren-answer-environment)))
+                         (let ((E (eval `(lambda ,(cadr expr) ,(last expr)) (μkanren-extension/environment)))
\end{scheme}

\subsection{Two primitives: forget and post-unify}\label{sec:design-forget}

The protocol exports two procedures besides the registry.

\scm{(μkanren-state-forget α s)} unbinds $α$:
\begin{scheme}
+  (define (μkanren-state-forget α s)
+    (let1 (S (μkanren-state-S s))
+      (μkanren-state-with s S: (update/sbral (μkanren-var-index/sbral S α) *μkanren-unbound* S))))
\end{scheme}
A \texttt{post} hook uses it to take back a binding that \scm{=°} has just made, turning it into a
constraint of its own, and leaves that binding out of the bindings it answers. Since \texttt{post} runs
before D, A and T are verified, those stores never see the binding. The λ layer needs this when the scope
check of $X:=t$ cannot be decided yet, because $t$ contains a non-pattern that only a later binding can
resolve: the binding becomes the suspended equation $X\doteq t$, as Teyjus would delay it. For instance
$\exists q\,\exists G\,\exists H\,\forall y.\ q = G\,(H\,y)$ leaves \scm{q} unbound with a residual; a later
$G:=\abs{a}{\cst k}$ solves it with $q=\cst k$ (\cref{ch:module}).

\scm{μkanren-post-unify}, shown above, is exported for goals of an extension that unify on their own
and must then run the same chain as \scm{=°}. The λ layer calls it at ∀ exit, where the raising of the
variables of a residual constraint makes new bindings that must wake the suspended equations and be
checked against D, A and T.

\subsection{One generic change to absent°}\label{sec:design-absent}

The third review found a defect of micro itself, independent of λ-terms. Micro decided an atom against
an \emph{absento} datum with \scm{equal?} only, although a datum that holds a variable may still become
that atom. With the unedited micro:
\begin{scheme}
(pretty-print (°->list #f (fresh° (q D) (absent° D q) (=° q 'a) (=° D 'a))))
(pretty-print (°->list #f (fresh° (q D) (=° q 'a) (=° D 'a) (absent° D q))))
(pretty-print (°->list #f (fresh° (q D) (absent° D q) (=° q 'a))))
\end{scheme}
\begin{answers}
((λ () 'a))
((λ () 'a))
((λ () 'a))
\end{answers}
The first two answers claim that \scm{a} is absent from \scm{a}. With the edited micro, which decides
such an atom by \scm{μkanren-tag-unify} (as it already did for a compound) and watches it while it is
pending:
\begin{answers}
()
()
((λ () 'a))
\end{answers}
The third answer stays, correctly: \scm{D} is not in the answer and can be chosen different from
\scm{a}. Two smaller generic changes come with this one. An answer keeps a pending atom constraint when
its datum names a variable of the answer, and prints that datum as a term,
\scm{(absent? (@ α 'b) 'a)}. And when a type turns an \emph{absento} into a disequality prefix of D (micro's
subsumption), a datum that holds an extension term stays in T instead (the predicate
\scm{μkanren-extension/holds-term?}), so D never holds one:
\begin{scheme}
+      ((((,α* . ,tag) . ,T*) ⊣ (and (equal? α α*) (not (μkanren-extension/holds-term? (μkanren-tag-name tag)))))
+        (let1 (D* (μkanren-ext-D α tag D s)) (μkanren-update-D/T+ α T+ D* T* s)))
\end{scheme}

\subsection{The two inherited invariants}\label{sec:design-invariants}

A client of the protocol must respect two invariants that micro already relies on.

\begin{definition}[Protocol invariants]\label{def:design-invariants}\leavevmode
\begin{enumerate}[label=(I\arabic*)]
\item \textbf{eq?-entailment.} A hook that changes nothing returns the very state object it was given.
\item \textbf{Extended bindings.} A hook that binds returns the list of bindings it received extended
  with every binding it made, pruning and lowering included.
\end{enumerate}
\end{definition}

(I1) is \cref{prop:micro-eq} carried over to the hooks. Micro reads \enquote{unification gave back $s$}
as \enquote{the terms are already equal}: \scm{≠°} then fails, \scm{=°} skips re-verification, and the
re-verification of a D prefix declares it violated. A hook that rebuilt an unchanged state would make
\scm{(≠° (ƛ (a) a) (ƛ (b) b))} look satisfiable (the λ analogue of \scm{(≠° 1 1)} in the pitfall
after \cref{prop:micro-eq}). For the λ
layer this has a concrete consequence: the names that the ξ rule opens under binders live outside the
state, so a ξ step that binds nothing returns $s$ itself. So the λ disequality below is recognised
as an equality that already holds, and fails, exactly as micro's \scm{(≠° 1 1)} does:
\begin{scheme}
(pretty-print (°->list #f (fresh° (q) (≠° (ƛ (a) a) (ƛ (b) b)))))
(pretty-print (°->list #f (fresh° (q) (≠° 1 1))))
\end{scheme}
\begin{answers}
()
()
\end{answers}

(I2) is what keeps the stores honest. \scm{μkanren-verify-D/touched} re-verifies only the prefixes
whose variables appear among the new bindings; a binding made by a hook and not reported would leave a
violated disequality in D. This is Trick~1's lesson (\cref{sec:micro-noedit}): a state produced by raw
unification must go through the D/A/T chain with its exact bindings.

\begin{pitfall}[A hook that forgets a binding must also drop it from the bindings]
(I2) has a converse. When \texttt{post} takes a binding back with \scm{μkanren-state-forget}, it must
remove it from the bindings it answers; otherwise the stores would be re-verified against a binding that
S no longer holds. The λ layer's \scm{μkanren-enforce} does both, through \scm{μkanren-drop}.
\end{pitfall}

\subsection{From closed points to hooks}\label{sec:design-map}

\Cref{tab:design-map} reads the protocol against \cref{tab:micro-open}: for each closed point of
\micro, what the edit opens and what it leaves closed on purpose.

\begin{table}[ht]
\centering\footnotesize
\caption{Closed points of \micro{} and the hooks that open them.}
\label{tab:design-map}
\renewcommand{\arraystretch}{1.2}
\begin{tabularx}{\textwidth}{@{}>{\raggedright\arraybackslash}p{3.8cm}>{\raggedright\arraybackslash}X>{\raggedright\arraybackslash}p{3.6cm}@{}}
\toprule
closed point (\cref{tab:micro-open}) & opened by & still closed\\
\midrule
the field list of the state, its positional constructions & the field \texttt{X}, copied at the two constructions & the six-ary \texttt{μkanren-state-match}\\
the unification \texttt{cond} & \texttt{unify} in the structure--structure and \texttt{else} clauses & the variable clauses, pair--pair\\
the binding site \texttt{μkanren-state-update} & \texttt{atom?} inside the occurs check; \texttt{occurs} on the failure path; the scope check moved to \texttt{post} & var--var orientation (older to newer)\\
the post-\texttt{=°} chain & \texttt{post}, threaded; \texttt{μkanren-post-unify} exported; \texttt{μkanren-state-forget} & ---\\
reification (D/A/T only) & \texttt{value}, \texttt{repr}, \texttt{terms}, \texttt{forms}; \texttt{parts} in projection mode & purification and subsumption of micro's stores\\
\texttt{°->§}: fixed empty initial state & not needed: \texttt{X = '()} means empty; \texttt{environment} for ground answers & fixed query index 0\\
\midrule
(listed as open: \texttt{≠°}, \texttt{absent°} and types work on any record, slot by slot) & \texttt{diseq}, \texttt{parts}, \texttt{tag} & ---\\
\bottomrule
\end{tabularx}
\end{table}

The three needs listed at the end of \cref{ch:micro} correspond to three rows: \enquote{a named place in
the state} is \texttt{X}, \enquote{a call after every binding} is \texttt{post}, and \enquote{a say in
reification} is \texttt{value}, \texttt{repr}, \texttt{terms} and \texttt{forms}. The rest of the
protocol exists because micro's stores and occurs check, written for first-order data, look into records
slot by slot.

\begin{tutorial}[One \texttt{=°} on a λ-term, through micro and the hooks]
To watch the protocol at work we re-register the λ extension under its own name with some hooks wrapped
by a printer. Registering under the same name replaces it (\cref{sec:design-registry}). Micro exports no
accessor of the hooks, so the script reads the slots of the registered record directly, with
\scm{record-instance-slot} of \scm{(chicken memory representation)} (imported with \scm{srfi-1},
besides \scm{(aux base)}, \scm{(aux stream)}, \micro{} and \scm{(aux kanren lambda)}); this is for tracing only.
\begin{scheme}
(define fields '(name term? atom? unify occurs post diseq tag parts
                 value repr terms forms environment))
(define (orig h)
  (record-instance-slot μkanren-λ-extension (list-index (λ (f) (eq? f h)) fields)))
(define (D t) (if (μkanren-λ-term? t) (μkanren-datum t) t))
(define (B as) (map (λ (p) (list (car p) ':= (D (cdr p)))) as))
(define (say . xs) (display "  ") (for-each display xs) (newline))
(define traced
  (make-μkanren-extension 'λProlog     ; same name: replaces the λ extension
    term?: (orig 'term?)
    atom?: (λ (u) (let1 (r ((orig 'atom?) u))
                    (say "atom?  " (D u) " => " r) r))
    unify: (λ (u v s as)
             (say "unify  " (D u) " =? " (D v))
             (let1 (r ((orig 'unify) u v s as))
               (say "unify  => " (and r (B (cadr r)))) r))
    occurs: (orig 'occurs)
    post: (λ (s as)
            (say "post   " (B as) " at depth " (μkanren-λ-depth s))
            (let1 (r ((orig 'post) s as))
              (say "post   => " (and r (B (cadr r)))) r))
    diseq: (orig 'diseq) tag: (orig 'tag) parts: (orig 'parts)
    value: (orig 'value) repr: (orig 'repr)
    terms: (λ (s) (let1 (r ((orig 'terms) s)) (say "terms  => " r) r))
    forms: (λ (s s*) (let1 (r ((orig 'forms) s s*)) (say "forms  => " r) r))
    environment: (orig 'environment)))
(μkanren-extension-register! traced)
(pretty-print (°->list/datum (fresh° (q) (∀° (x) (=° (@ q x) `(g ,x ,x))))))
(pretty-print (°->list #f (fresh° (q) (∀° (x) (=° q x)))))
\end{scheme}
\begin{answers}
  unify  (@ _0 (eigen x)) =? (g (eigen x) (eigen x))
  atom?  (ƛ (x) (g x x)) => #f
  atom?  (bvar 0) => #t
  atom?  (bvar 0) => #t
  unify  => ((_0 := (ƛ (x) (g x x))))
  post   ((_0 := (ƛ (x) (g x x)))) at depth 1
  post   => ((_0 := (ƛ (x) (g x x))))
  terms  => ()
  forms  => ()
((ƛ (x) (g x x)))
  atom?  (eigen x) => #t
  post   ((_0 := (eigen x))) at depth 1
  post   => #f
()
\end{answers}
\textbf{The first query}, \lp{pi x\ F x = g x x}, step by step.
\begin{enumerate}
\item \scm{∀°} raises the depth to 1 and makes the eigenvariable \scm{x} of level 1; \scm{q} is
  \scm{_0}, of level 0.
\item Micro's \scm{=°} calls \scm{μkanren-state-unify}. Both sides are walked: the left one is an
  application record, the right one a pair. No variable clause applies, the left side is a structure
  but the right one is not, and the pair--pair clause does not apply: the final \scm{else} finds the
  owning extension and calls \texttt{unify}.
\item The hook head-normalises both sides. The left one is flexible, \scm{_0} applied to the pattern
  \scm{(x)}: one eigenvariable, whose level 1 exceeds the level 0 of \scm{q}. So this is flex--rigid. Variable
  elimination replaces \scm{x} by the bound variable of the new abstraction, and the hook calls micro
  back to bind \scm{_0 := (ƛ (x) (g x x))}.
\item That binding goes through micro's binding site and its occurs check. The check looks into the
  abstraction (\texttt{atom?} is false) and stops at each \scm{(bvar 0)} (\texttt{atom?} is true). No
  occurrence of \scm{_0}: the binding is made, and the hook returns it.
\item Back in \scm{=°}, the state has changed, so \scm{μkanren-post-unify} runs the \texttt{post} hook
  on the new binding. Levels exist, and \scm{_0} is below the highest level reached, so the scope check
  runs: the bound term holds no eigenvariable of level $>0$ (\scm{x} has become an index), and no
  variable to lower. The binding passes; E is empty, so nothing is woken. Then micro verifies D, A and
  T, all empty.
\item The \scm{∀°} body has answered; \scm{∀°} gives the state back depth 0 and discharges the
  constraints mentioning \scm{x} (there are none).
\item Reification asks the extension for the terms of its residual constraints (none) and for its
  forms (none); the \texttt{value} and \texttt{repr} hooks, not traced here, print the abstraction.
\end{enumerate}
\textbf{The second query}, \lp{sigma Y\ pi x\ Y = x}, never calls \texttt{unify}: \scm{q} is a variable,
so micro's variable clause binds it directly. The occurs check meets the eigenvariable, an atom, and
does not look into it. Micro makes the binding \scm{_0 := x}; but the \texttt{post} hook sees that
\scm{x} has level 1 and \scm{_0} level 0, which violates condition (2) of A1, and answers \scm{#f}.
\scm{=°} fails before any other goal could see the binding. This is A1(c) at work: the binding was
made by micro's own clause, and the extension still had the last word.
\end{tutorial}

\subsection{Cost}\label{sec:design-cost}

\Cref{tab:roadmap-perf-protocol} in \cref{ch:roadmap} gives the cost of the protocol on first-order
search, compiled with \texttt{csc -O2 -d1}: the unedited micro against the edited micro, without and
with \scm{(aux kanren lambda)} loaded but unused.
The protocol costs at most 1\% end to end on \scm{append°} (5923 to 5959~ms in round 2), within the
2.5\% round-to-round noise of the unedited binary. With the λ extension loaded, the disequality-heavy
perm6 costs 8--12\% (d2 cost 45--50\%): the remainder is the \texttt{post} hook's look at the new
bindings while D is not empty. One cost is not removed and is documented: once any \scm{∀°} has run,
levels exist, and every later binding below the top level gets the scope check, which costs first-order
search 1.7 to 2.3~times (\cref{ch:roadmap} plans one walk for the occurs and the scope check).

% ================================================================================================
\section{Review history}\label{sec:design-reviews}

After the choice of d2 and the grafts, the proposal went through three adversarial review rounds. Each
round attacked the proposal itself with hand-written probes and with randomly generated goals, run
with their conjunctions in both orders. Every report was first reproduced, then fixed, a code fix
with a regression case (in round 3 each new case was checked to fail on the proposal before the fix).

\begin{table}[ht]
\centering\footnotesize
\caption{The three adversarial review rounds (design notes, §10.4--10.6).}
\label{tab:design-reviews}
\renewcommand{\arraystretch}{1.2}
\begin{tabularx}{\textwidth}{@{}lrr>{\raggedright\arraybackslash}X@{}}
\toprule
round & reports & regression cases & typical defects\\
\midrule
1 & 13 & 12 (44 assertions) & disequality prefixes holding an eigenvariable left in D (a free \texttt{(eigen x)} printed, a crash); pruning inside an argument that a non-pattern may drop; occurs check not modulo β; raw record slots in projected \emph{absento}; a binder hint capturing \texttt{cons} in ground answers\\
2 & 8 & 6 (31 assertions) & raising at ∀ exit not closing every constraint; η-redexes created by later bindings; occurs failure on non-rigid arguments; the extensional reading absorbed by T; first-order cost after a \texttt{∀°}\\
3 & 12 & 9 (53 assertions) & non-patterns with functional arguments; disequalities over unreachable arguments; \emph{absento} of an eigenvariable under a type; user types on λ-terms; η-short occurs check; a datum reducing to an atom; two types on one flexible term; first-order cost with D or A non-empty; λ-terms in the show module's cells view\\
\bottomrule
\end{tabularx}
\end{table}

All 33 reports were real. In round 1 the twelve code defects were fixed, each pinned by a case, and one
documentation error was corrected. In round 2 two reports were one defect; six were fixed with
regression cases, and the seventh, the cost of first-order search after a \scm{∀°}, was measured and
documented rather than changed, since the cheap fix would give up the guarantee that host code cannot
smuggle an eigenvariable into an older variable. In round 3 three pairs of reports were one defect each;
all nine defects were fixed with a regression case each. Round 3 also produced the generic change to
\scm{absent°} of \cref{sec:design-absent}. The defects themselves belong to the λ layer and are discussed
where it is described (\cref{ch:module}).

% ================================================================================================
\section{Summary}\label{sec:design-summary}

\begin{itemize}
\item The requirements are nine acceptance items: the scope of ∀ under Nadathur, Jayaraman and Kwon's
  binding conditions, for \emph{every} binding (A1); \rulename{augment} with the book's clause order and variable
  sharing (A2); type inference as the test of ∀ and ⊃ together (A3); λ-tree syntax with α, β and η (A4);
  $\Llam$ unification with mgus and delayed non-patterns (A5); λProlog's answers where Ma and Friedman
  differ (A6); the book's programs (A7); micro's stores on λ-terms (A8); micro's suites unchanged
  (A9). Order independence runs through all of them. They are pinned by 83 cases and 320 assertions.
\item Of the four meanings of \scm{=°} on λ-terms, pattern unification with universe levels is
  retained: it is λProlog's, it is decidable and unitary (one state per \scm{=°}), and levels give
  \rulename{generic} for one integer per variable and one comparison per binding.
\item Three prototypes were built: a layer with no edit (d1), a hooks protocol (d2), a native engine
  (d3). d1 cannot check bindings made by micro's \scm{=°}, so eigenvariables leak through \scm{cons°};
  d3 had a crash and an η bug that made \scm{≠°} unsound; d2 had seven non-soundness defects. The
  reviews scored them 15.5, 23 and 21; d2 became the base, and d3's pedagogical strengths and d1's
  precision were grafted onto it.
\item The protocol adds a seventh state field \texttt{X} that micro copies and never reads, a global
  registry of extensions dispatched by \texttt{term?}, and thirteen hooks: \texttt{atom?} and
  \texttt{occurs} in the binding site, \texttt{unify} in the unification \scm{cond}, \texttt{post} in the
  chain after \scm{=°}, \texttt{diseq}, \texttt{tag} and \texttt{parts} in the stores, \texttt{value},
  \texttt{repr}, \texttt{terms} and \texttt{forms} in reification, \texttt{environment} in the runner.
\item The scope check lives in \texttt{post}: every binding passes through some \scm{=°}, and every
  \scm{=°} ends in \scm{μkanren-post-unify}, which runs the hooks before D, A and T are verified.
  An extension goal that unifies on its own must end in \scm{μkanren-post-unify} too.
  \scm{μkanren-state-forget} lets \texttt{post} turn a binding into a suspended equation.
\item A client must return the very state when it changes nothing (eq?-entailment) and report every
  binding it makes (extended bindings).
\item The edit is +235/−41 lines, λ-agnostic, and costs about 1\% on first-order search; micro's 125
  cases pass unchanged with and without the λ layer. Three adversarial rounds sent 13, 8 and 12 reports,
  all real; every defect but one documented cost was fixed with a regression case.
\end{itemize}
